This thesis explores the application of deep learning, specifically fully connected neural networks (FCNN), to predict the hardening behavior of lithium-ion pouch cells under mechanical loading. Traditional two-stage analytical models often struggle to capture complex deformation responses, particularly in varied indentation scenarios. By integrating a neural network-based hardening law into finite element simulations, this study aims to improve the prediction of yield stress from plastic strain and enhance force-displacement matching.

This thesis introduces various feature selection methods, focusing on the gradient of the loss function. Various model configurations, machine learning hyperparameters, and selection methods are evaluated to optimize performance. The results demonstrate that while deep learning can effectively model and calibrate for a von Mises yield surface and its associated hardening response, there remains significant room for improvement in achieving consistent accuracy and stability.

